\documentclass[
	fontsize=12pt,
	paper=a4,
	numbers=noenddot,
	parskip=half,
	openany
]{scrbook}
\usepackage{tcolorbox}

\include{../../../for_latex/preamble}
% ##################################################
% Gemeinsames Setup fuer die Probeklausuren
% (Java-Listings, Hinweis-/Beispiel-Boxen, Struktogramm-Stile in TikZ)
% wird nach \include{../for_latex/preamble} per \input eingebunden
% ##################################################

% --- Java Code-Highlighting ---
\definecolor{codegray}{rgb}{0.5,0.5,0.5}
\definecolor{codepurple}{rgb}{0.58,0,0.82}
\definecolor{backcolour}{rgb}{0.95,0.95,0.92}
\definecolor{keywordcolor}{rgb}{0.8,0.47,0.13}

\lstdefinestyle{javastyle}{
	backgroundcolor=\color{backcolour},
	commentstyle=\color{codegray},
	keywordstyle=\color{keywordcolor}\bfseries,
	numberstyle=\tiny\color{codegray},
	stringstyle=\color{codepurple},
	basicstyle=\ttfamily\footnotesize,
	breakatwhitespace=false,
	breaklines=true,
	captionpos=b,
	keepspaces=true,
	numbers=left,
	numbersep=5pt,
	showspaces=false,
	showstringspaces=false,
	showtabs=false,
	tabsize=2,
	language=Java,
	frame=single,
	rulecolor=\color{black!30}
}

\lstset{style=javastyle}

% --- Hinweis-/Beispiel-/Formel-Umgebungen (schlicht, ohne farbige Box) ---
\newenvironment{hinweis}[1][Hinweis]%
	{\par\smallskip\noindent\textbf{#1.}\enspace\ignorespaces}%
	{\par\smallskip}

\newenvironment{beispiel}[1][Beispiel]%
	{\par\smallskip\noindent\textbf{#1:}\par\nobreak\vspace{2pt}}%
	{\par\smallskip}

\newenvironment{formeln}[1][Gegebene Formeln]%
	{\par\smallskip\noindent\textbf{#1:}\par\nobreak\vspace{2pt}}%
	{\par\smallskip}

% ##################################################
% Struktogramm (Nassi-Shneiderman) in reinem TikZ
% ##################################################
\newlength{\nswidth}\setlength{\nswidth}{11.5cm}   % Gesamtbreite
\newlength{\nsind}\setlength{\nsind}{0.55cm}        % Einrueckung je Schleifenebene
% vorberechnete Versatzbreiten fuer Verzweigungen (in Koordinaten nutzbar)
\newlength{\nshalf}\setlength{\nshalf}{\dimexpr\nswidth/2\relax}            % halbe Gesamtbreite
\newlength{\nshalfi}\setlength{\nshalfi}{\dimexpr(\nswidth-\nsind)/2\relax}  % halbe Breite, 1x eingerueckt
\newlength{\nshalfii}\setlength{\nshalfii}{\dimexpr(\nswidth-2\nsind)/2\relax} % halbe Breite, 2x eingerueckt

\tikzset{
	% Basis fuer Anweisungs-/Verzweigungsbloecke
	nb/.style={draw, line width=0.4pt, anchor=north west, inner sep=4pt,
		align=left, font=\footnotesize, fill=white, minimum height=0.8cm},
	% volle Breite
	nbfull/.style={nb, minimum width=\nswidth,
		text width=\dimexpr\nswidth-8pt\relax},
	% eine Ebene eingerueckt (Schleifenrumpf)
	nbi/.style={nb, minimum width=\dimexpr\nswidth-\nsind\relax,
		text width=\dimexpr\nswidth-\nsind-8pt\relax},
	% zwei Ebenen eingerueckt (Schleife in Schleife)
	nbii/.style={nb, minimum width=\dimexpr\nswidth-2\nsind\relax,
		text width=\dimexpr\nswidth-2\nsind-8pt\relax},
	% Verzweigung: halbe Breiten
	nbh/.style={nb, minimum width=\dimexpr\nswidth/2\relax,
		text width=\dimexpr\nswidth/2-8pt\relax},
	nbih/.style={nb, minimum width=\dimexpr(\nswidth-\nsind)/2\relax,
		text width=\dimexpr(\nswidth-\nsind)/2-8pt\relax},
	% Verzweigung halbe Breite, zwei Ebenen eingerueckt
	nbiih/.style={nb, minimum width=\nshalfii,
		text width=\dimexpr\nshalfii-8pt\relax},
	% Kopf einer Verzweigung (das Dreieck), inner sep 0
	nc/.style={draw, line width=0.4pt, anchor=north west, inner sep=0pt,
		minimum height=1.0cm, fill=white},
	ncfull/.style={nc, minimum width=\nswidth},
	nci/.style={nc, minimum width=\dimexpr\nswidth-\nsind\relax},
	ncii/.style={nc, minimum width=\dimexpr\nswidth-2\nsind\relax},
}

% Zeichnet das Dreieck + Beschriftungen in einen bereits gesetzten nc-Knoten.
% #1 Knotenname  #2 Bedingung  #3 ja-Label  #4 nein-Label
\newcommand{\nscond}[4]{%
	\draw[line width=0.4pt] (#1.north west) -- (#1.south);
	\draw[line width=0.4pt] (#1.north east) -- (#1.south);
	\node[font=\scriptsize, anchor=north, inner sep=1pt] at ($(#1.north)+(0,-1.5pt)$) {#2};
	\node[font=\scriptsize, anchor=south west, inner sep=1pt] at ($(#1.south west)+(2.5pt,1.5pt)$) {#3};
	\node[font=\scriptsize, anchor=south east, inner sep=1pt] at ($(#1.south east)+(-2.5pt,1.5pt)$) {#4};
}

% Zeichnet den L-Rahmen einer Schleife (linker + unterer Rand der Einrueckspalte).
% #1 Kopfknoten der Schleife  #2 letzter Rumpfknoten
\newcommand{\nsframe}[2]{%
	\draw[line width=0.4pt] (#1.south west) |- (#2.south west);
}


\title{Probeklausur 3}
\author{Luis Staudt}
\date{}

\begin{document}

% --- Titelblock ---
\begin{center}
	\rule{\textwidth}{0.8pt}\\[0.6em]
	{\LARGE\bfseries Probeklausur 3}\\[0.3em]
	{\large Grundlagen der Programmierung}\\[0.7em]
	\begin{tabular}{r l @{\hspace{1.4cm}} r l}
		\textbf{Studiengänge:}    & PEB \& MVB            & \textbf{Autor:}        & Luis Staudt \\
		\textbf{Semester:}        & Sommersemester 2026   & \textbf{Schwierigkeit:}& mittel \\
		\textbf{Bearbeitungszeit:}& 90 Minuten           & \textbf{Gesamtpunkte:} & 90 \\
	\end{tabular}\\[0.7em]
	\rule{\textwidth}{0.8pt}
\end{center}
\vspace{0.5em}

\noindent\textit{Bearbeiten Sie alle fünf Aufgaben. Achten Sie auf sauberen, kompilierfähigen Java-Quelltext. Hilfsmittel: OpenBook, keine elektronischen Geräte.}

\pagestyle{fancy}
\fancyhf{}
\lhead{\small Probeklausur 3, Grundlagen der Programmierung}
\rhead{\small SoSe 2026 (mittel)}
\cfoot{\thepage}
\renewcommand{\headrulewidth}{0.4pt}

% =============================================
\clearpage
\section*{Aufgabe 1: Struktogramm in Java umsetzen \hfill (24 Punkte)}

Das folgende Struktogramm liest Messwerte ein, zählt die gültigen Werte (innerhalb eines Bereichs) und ermittelt deren Mittelwert sowie Minimum und Maximum.
Anschließend sucht es den \emph{ersten} Ausreißer (Abweichung vom Mittelwert über einer Grenze).

Setzen Sie das Struktogramm \textbf{vollständig in ein lauffähiges Java-Programm} um.

\begin{center}
\begin{tikzpicture}
	\node[nbfull] (a1) at (0,0) {lies $n$, \texttt{unten}, \texttt{oben} und \texttt{grenze} ein};
	\node[nbfull] (a2) at (a1.south west) {lege Feld \texttt{werte} mit $n$ Plätzen an};
	\node[nbfull] (h1) at (a2.south west) {\textbf{für} $i = 0,\,1,\,\ldots,\,n-1$};
	\node[nbi]    (b1) at ($(h1.south west)+(\nsind,0)$) {lies \texttt{werte[i]} ein};
	\nsframe{h1}{b1}
	\node[nbfull] (a3) at ($(b1.south west)+(-\nsind,0)$) {\texttt{gueltig=0; tsum=0; min=werte[0]; max=werte[0];}};

	% Schleife 2
	\node[nbfull] (h2) at (a3.south west) {\textbf{für} $i = 0,\,1,\,\ldots,\,n-1$};
	\node[nci]    (c1) at ($(h2.south west)+(\nsind,0)$) {};
	\nscond{c1}{\texttt{unten <= werte[i] <= oben}}{ja}{nein}
	\node[nbih]   (c1l) at (c1.south west) {\texttt{gueltig++; tsum += werte[i];}};
	\node[nbih]   (c1r) at ($(c1.south west)+(\nshalfi,0)$) {\textit{(nichts)}};
	\node[nci]    (c2) at (c1l.south west) {};
	\nscond{c2}{\texttt{werte[i] < min}}{ja}{nein}
	\node[nbih]   (c2l) at (c2.south west) {\texttt{min = werte[i];}};
	\node[nbih]   (c2r) at ($(c2.south west)+(\nshalfi,0)$) {\textit{(nichts)}};
	\node[nci]    (c3) at (c2l.south west) {};
	\nscond{c3}{\texttt{werte[i] > max}}{ja}{nein}
	\node[nbih]   (c3l) at (c3.south west) {\texttt{max = werte[i];}};
	\node[nbih]   (c3r) at ($(c3.south west)+(\nshalfi,0)$) {\textit{(nichts)}};
	\nsframe{h2}{c3l}

	% Mittelwert (Verzweigung)
	\node[ncfull] (m1) at ($(c3l.south west)+(-\nsind,0)$) {};
	\nscond{m1}{\texttt{gueltig > 0}}{ja}{nein}
	\node[nbh]    (m1l) at (m1.south west) {\texttt{mittel = tsum / gueltig;}};
	\node[nbh]    (m1r) at ($(m1.south west)+(\nshalf,0)$) {\texttt{mittel = 0;}};
	\node[nbfull] (a4) at (m1l.south west) {\texttt{i = 0;}\qquad \texttt{gefunden = false;}};

	% Schleife 3 (while mit Flag)
	\node[nbfull] (h3) at (a4.south west) {\textbf{solange} (\texttt{nicht gefunden}) \textbf{und} (\texttt{i < n})};
	\node[nci]    (c4) at ($(h3.south west)+(\nsind,0)$) {};
	\nscond{c4}{$|\,\mathtt{werte[i]}-\mathtt{mittel}\,| > \mathtt{grenze}$}{ja}{nein}
	\node[nbih]   (c4l) at (c4.south west) {\texttt{gefunden = true; pos = i;}};
	\node[nbih]   (c4r) at ($(c4.south west)+(\nshalfi,0)$) {\textit{(nichts)}};
	\node[nbi]    (p1) at (c4l.south west) {\texttt{i = i + 1;}};
	\nsframe{h3}{p1}

	% Ausgabe (Verzweigung)
	\node[ncfull] (e1) at ($(p1.south west)+(-\nsind,0)$) {};
	\nscond{e1}{\texttt{gefunden == true}}{ja}{nein}
	\node[nbh]    (e1l) at (e1.south west) {drucke \texttt{pos}};
	\node[nbh]    (e1r) at ($(e1.south west)+(\nshalf,0)$) {drucke: kein Ausreißer};
\end{tikzpicture}
\end{center}

% =============================================
\clearpage
\section*{Aufgabe 2: Widerstandsschaltung auswählen \hfill (14 Punkte)}

Zwei oder drei Widerstände werden zusammengeschaltet, anschließend wird der Strom bei gegebener Spannung~$U$ berechnet.

\begin{formeln}
\[
	\text{(1) Reihe:}\quad R = R_1 + R_2 + R_3
	\qquad
	\text{(2) Parallel (zwei):}\quad R = \frac{R_1 \cdot R_2}{R_1 + R_2}
\]
\[
	\text{danach in beiden Fällen:}\quad I = \frac{U}{R}
\]
\end{formeln}

Schreiben Sie ein Programm, das die Schaltungsart per Zahl (\texttt{1} oder \texttt{2}) wählen lässt, die Widerstände und die Spannung~$U$ abfragt, den Gesamtwiderstand~$R$ mit der \textbf{richtigen} Formel und anschließend den Strom~$I$ berechnet. Fangen Sie den Sonderfall $R = 0$ (Kurzschluss) ab und geben Sie dann eine entsprechende Meldung statt einer Division aus.

\begin{beispiel}[Beispieldialog]
\begin{lstlisting}[language={}]
Schaltung (1=Reihe, 2=Parallel): 2
R1 [Ohm]: 100
R2 [Ohm]: 100
Spannung U [V]: 10
Gesamtwiderstand R = 50.0 Ohm
Strom I = 0.2 A
\end{lstlisting}
\end{beispiel}

% =============================================
\clearpage
\section*{Aufgabe 3: Klasse \texttt{Kreis} \hfill (20 Punkte)}

Gegeben ist die Klasse \texttt{Punkt} (Attribute \texttt{xKoordinate}, \texttt{yKoordinate} als \texttt{double}).
Implementieren Sie eine Klasse \texttt{Kreis}, die einen Mittelpunkt (ein \texttt{Punkt}) und einen Radius besitzt.
Das Gerüst ist vorgegeben. Ergänzen Sie:

\begin{itemize}[nosep]
	\item einen \textbf{Konstruktor} mit Mittelpunkt und Radius. Ist der Radius $\le 0$, wird stattdessen der Radius~\texttt{1} verwendet.
	\item die Methode \texttt{berechneFlaeche()} ($A = \pi r^2$).
	\item die Methode \texttt{enthaeltPunkt(Punkt p)}, die \texttt{true} zurückgibt, wenn \texttt{p} innerhalb des Kreises oder auf dem Rand liegt. \emph{Verwenden Sie dazu \texttt{berechneAbstandZuPunkt}.}
\end{itemize}

\begin{lstlisting}
public class Kreis {
    public Punkt mittelpunkt;  // Daten
    public double radius;      // Daten
    // Konstruktor und Methoden ergaenzen
}
\end{lstlisting}

\begin{lstlisting}
Punkt mp = new Punkt();
mp.xKoordinate = 2;  mp.yKoordinate = 3;
Kreis k = new Kreis(mp, 5);
System.out.println("Flaeche: " + k.berechneFlaeche());

Punkt p = new Punkt();
p.xKoordinate = 4;  p.yKoordinate = 4;
if (k.enthaeltPunkt(p))
    System.out.println("p liegt im Kreis");
\end{lstlisting}

% =============================================
\clearpage
\section*{Aufgabe 4: Fachwerk, Mittelpunkt des höchstbelasteten Stabes \hfill (16 Punkte)}

Gegeben ist die aus Vorlesung und Übung bekannte Datei \texttt{fachwerk.jar}.
Für ein gewähltes Fachwerk soll nach der Berechnung der Stab mit der betragsmäßig \textbf{größten Stabkraft} bestimmt werden.
Geben Sie dessen Stabnummer, die Kraftart (Zug/Druck) sowie die Koordinaten seines \textbf{Mittelpunkts} aus.

\begin{formeln}[Mittelpunkt eines Stabes]
\[ x_m = \frac{x_a + x_e}{2} \qquad y_m = \frac{y_a + y_e}{2} \]
wobei $a$ der Anfangs- und $e$ der Endknoten des Stabes ist.
\end{formeln}

\begin{center}
\begin{tikzpicture}[>=Stealth, scale=1.0,
	knoten/.style={circle, fill=black, inner sep=1.6pt},
	stab/.style={thick, gray!70!black}]
	\coordinate (A) at (0,0); \coordinate (B) at (2.2,0); \coordinate (C) at (4.4,0); \coordinate (D) at (6.6,0);
	\coordinate (E) at (1.1,1.8); \coordinate (F) at (3.3,1.8); \coordinate (G) at (5.5,1.8);
	\draw[stab] (A)--(B); \draw[stab] (B)--(C); \draw[stab] (C)--(D);
	\draw[stab] (A)--(E); \draw[stab] (E)--(B); \draw[stab] (E)--(F);
	\draw[stab] (F)--(B); \draw[stab] (F)--(C); \draw[stab] (F)--(G);
	\draw[stab] (G)--(C); \draw[stab] (G)--(D);
	\foreach \p in {A,B,C,D,E,F,G} \node[knoten] at (\p) {};
	\draw[->, thick] (B)--++(0,-0.9); \draw[->, thick] (C)--++(0,-0.9);
\end{tikzpicture}\\[-0.3em]
{\footnotesize\itshape Schematische Darstellung eines Fachwerks.}
\end{center}

\begin{hinweis}[Benötigte Methoden aus \texttt{fachwerk.jar}]
	\texttt{gibAnzahlStaebe()}, \texttt{gibStab(i)},\\ \texttt{gibStabkraft()}, \texttt{gibKnotenNrA()}, \texttt{gibKnotenNrE()}, \texttt{gibKnoten(i)}, \texttt{gibX()}, \texttt{gibY()}.
\end{hinweis}

\begin{lstlisting}
public class Fachwerkaufgabe {
    public static void main(String[] args) {
        System.out.println("Welches Fachwerk wollen Sie rechnen (0-4)?");
        int wahl = Tastatureingabe.leseInt();
        Fachwerk fw = new Fachwerk(wahl);
        fw.berechneGroessen();
        // hier erweitern: hoechstbelasteten Stab finden, Mittelpunkt ausgeben
    }
}
\end{lstlisting}

\begin{beispiel}[Mögliche Ausgabe]
\begin{lstlisting}[language={}]
Hoechstbelasteter Stab: Stab 3 (Druckkraft), Kraft = -62.5
Mittelpunkt: x = 3.0   y = 1.5
\end{lstlisting}
\end{beispiel}

% =============================================
\clearpage
\section*{Aufgabe 5: Entladekurve in eine Datei schreiben \hfill (16 Punkte)}

Ein Kondensator entlädt sich; die Spannung folgt
\[ U(t) = U_0 \cdot e^{-t/\tau}, \qquad U_0 = 10~\text{V}, \quad \tau = 5~\text{s}. \]

Schreiben Sie ein Programm, das für $t = 0~\text{s}$ bis $t = 20~\text{s}$ in Schritten von $0{,}5~\text{s}$ jeweils $t$ und $U(t)$ in die Datei \texttt{entladung.txt} schreibt.
Markieren Sie zusätzlich in der jeweiligen Zeile, sobald $U(t) < 0{,}37 \cdot U_0$ unterschritten wird (z.\,B.\ mit dem Zusatz \texttt{*}).

\begin{hinweis}[Hinweis]
	Die Exponentialfunktion ist \texttt{Math.exp(...)}.
	Datei: \lstinline|new FileWriter("entladung.txt")|, je Zeile \texttt{write(...)} mit \lstinline|"\n"|, am Ende \texttt{close()}.
	Eine Kopfzeile ist erwünscht.
\end{hinweis}

\begin{beispiel}[Erwarteter Dateiinhalt (Auszug)]
\begin{lstlisting}[language={}]
t [s]    U [V]
0.0      10.0
0.5      9.048
...
5.0      3.679   *
...
20.0     0.183   *
\end{lstlisting}
\end{beispiel}

\end{document}
